\section*{Overview}

Sqrt decomposition is the middle ground between a raw array and a full segment tree. It groups the array into blocks
and stores one summary per block, which is often enough to answer queries within the required limits.

\section*{When to Use It}

Use it when:

\begin{itemize}
  \item the operation is simple enough to summarize per block,
  \item a segment tree would work but feels heavier than necessary,
  \item the constraints are moderate and a \(O(\sqrt{n})\) tradeoff is acceptable.
\end{itemize}

\section*{Core Idea}

Partition the array into blocks of size about \(\sqrt{n}\). For a range query:

\begin{itemize}
  \item scan the partial blocks at the ends directly,
  \item use the precomputed summaries for the fully covered middle blocks.
\end{itemize}

\section*{Key Insight}

Only \(O(\sqrt{n})\) blocks can be touched by one query. The only expensive work is on the two boundary blocks; the
middle is handled by block summaries.

\section*{Operations / Main Technique}

\begin{itemize}
  \item choose a block size,
  \item precompute one summary per block,
  \item answer queries by mixing direct scans and block summaries,
  \item update one position by fixing both the raw value and its block summary.
\end{itemize}

\paragraph{Worked example.}
For range sum with point updates, each block stores its total. A query over \([l, r]\) scans at most two partial blocks
and adds the totals of the whole blocks in between.

\section*{Correctness Intuition}

The range is partitioned into disjoint pieces: left partial block, some number of full blocks, and right partial block.
Each piece contributes exactly once, so the total is correct.

\section*{Complexity Analysis}

For block size \(B\):
\begin{itemize}
  \item point update: \(O(1)\),
  \item range query: \(O(B + n/B)\).
\end{itemize}
Choosing \(B \approx \sqrt{n}\) gives the usual \(O(\sqrt{n})\) bound.

\section*{Implementation}

The sample code implements point update and range-sum query. The same framework can support other mergeable block
statistics as long as a whole block can be summarized compactly.

\section*{Common Pitfalls}

\begin{itemize}
  \item Choosing a block size that is much too small or too large.
  \item Forgetting to update the block summary after a point update.
  \item Using sqrt decomposition when the query/update mix clearly calls for a stronger structure.
\end{itemize}

\section*{Variants / Extensions}

\begin{itemize}
  \item Range add with lazy block tags.
  \item Ordered statistics inside blocks after rebuilding.
  \item Mo's algorithm as a different block-based philosophy for offline queries.
\end{itemize}

\section*{Practice Problems}

\begin{itemize}
  \item Range sum with point updates.
  \item Range minimum with easy block recomputation.
  \item Frequency or order-statistic queries with sorted blocks.
\end{itemize}

\section*{References}

\begin{itemize}
  \item \href{https://cp-algorithms.com/data_structures/sqrt_decomposition.html}{cp-algorithms: Sqrt Decomposition}.
  \item \href{https://codeforces.com/catalog}{Codeforces Catalog}.
\end{itemize}
